\documentclass[border=15pt, multi, tikz]{standalone}
\usepackage{ctex}  % 引入ctex包以支持中文
\usepackage{import}
\subimport{../../layers/}{init}  % 引入layers初始化文件
\usetikzlibrary{positioning}
\usetikzlibrary{calc}

% 定义Transformer架构中使用的颜色
\definecolor{attention}{HTML}{F4D03F}     % 注意力机制 - 金黄色
\definecolor{feedforward}{HTML}{5DADE2}   % 前馈网络 - 蓝色
\definecolor{norm}{HTML}{D7BDE2}          % 归一化层 - 淡紫色
\definecolor{embedding}{HTML}{F1948A}     % 嵌入层 - 粉红色
\definecolor{linear}{HTML}{85C1E9}        % 线性层 - 浅蓝色
\definecolor{softmax}{HTML}{82E0AA}       % Softmax - 绿色
\definecolor{encoder}{HTML}{E8F6F3}       % 编码器背景 - 浅绿色
\definecolor{decoder}{HTML}{FDF2E9}       % 解码器背景 - 浅橙色

\begin{document}
\begin{tikzpicture}[
    % 定义连接线样式
    connect/.style={-latex, thick, black!70},
    skip/.style={-latex, thick, gray!60, dashed},
    % 定义文字样式
    layer_text/.style={font=\footnotesize\bfseries, text=black, align=center},
    title_text/.style={font=\large\bfseries, text=black}
]

%%%%%%%%%%%%%%%%%%%%%%%%%%%%%%%%%%%%%%%%%%%%%%%%%%%%%%%%%%%%%%%%%%%%%%%%%%%%%%%%%
% 编码器部分 (Encoder)
%%%%%%%%%%%%%%%%%%%%%%%%%%%%%%%%%%%%%%%%%%%%%%%%%%%%%%%%%%%%%%%%%%%%%%%%%%%%%%%%%

% 输入嵌入层 (Input Embedding)
\pic[shift={(0,0)}] {RoundedRectangle={
    name=input_embed,
    caption=,
    fill=embedding,
    width=3,
    height=0.8,
    cornerradius=5pt,
    bordercolor=black!50,
    borderwidth=0.8pt
}};
% 在矩形内部添加文字
\node[layer_text] at (input_embed-anchor) {输入嵌入};

% 位置编码 (Positional Encoding) - 使用圆形
\node[circle, draw=black!50, fill=gray!30, minimum size=0.8cm] (pos_enc_left) at (0, 1.5) {};
\node[layer_text] at (pos_enc_left) {PE};
\node[layer_text] at (-1, 1.5) {位置编码};

% 加法符号
\node[circle, draw=black!50, fill=white, minimum size=0.6cm] (add1) at (1.5, 1.5) {};
\node[layer_text] at (add1) {+};

% 编码器第一层
% Multi-Head Self-Attention
\pic[shift={(0,3)}] {RoundedRectangle={
    name=enc_attn1,
    caption=,
    fill=attention,
    width=3,
    height=0.8,
    cornerradius=5pt,
    bordercolor=black!50,
    borderwidth=0.8pt
}};
\node[layer_text] at (enc_attn1-anchor) {多头自注意力};

% Add & Norm 1
\pic[shift={(0,4.2)}] {RoundedRectangle={
    name=enc_norm1,
    caption=,
    fill=norm,
    width=3,
    height=0.6,
    cornerradius=5pt,
    bordercolor=black!50,
    borderwidth=0.8pt
}};
\node[layer_text, font=\scriptsize\bfseries] at (enc_norm1-anchor) {相加 \& 层归一化};

% Feed Forward Network
\pic[shift={(0,5.4)}] {RoundedRectangle={
    name=enc_ff1,
    caption=,
    fill=feedforward,
    width=3,
    height=0.8,
    cornerradius=5pt,
    bordercolor=black!50,
    borderwidth=0.8pt
}};
\node[layer_text] at (enc_ff1-anchor) {前馈神经网络};

% Add & Norm 2
\pic[shift={(0,6.6)}] {RoundedRectangle={
    name=enc_norm2,
    caption=,
    fill=norm,
    width=3,
    height=0.6,
    cornerradius=5pt,
    bordercolor=black!50,
    borderwidth=0.8pt
}};
\node[layer_text, font=\scriptsize\bfseries] at (enc_norm2-anchor) {相加 \& 层归一化};

% \draw[thick, encoder, rounded corners=10pt] (-0.2, 2.5) rectangle (3.2, 7.5);
% 编码器背景框 - 只框住4个主要组件，使用透明的编码器颜色填充
\draw[thick, fill=encoder, fill opacity=0.3, draw=encoder!80!black, rounded corners=10pt] (-0.2, 2.5) rectangle (3.2, 7.5);
\node[title_text, color=encoder!80!black, align=center] at (-1.5, 5.8) {编码器\\(Encoder)};

%%%%%%%%%%%%%%%%%%%%%%%%%%%%%%%%%%%%%%%%%%%%%%%%%%%%%%%%%%%%%%%%%%%%%%%%%%%%%%%%%
% 解码器部分 (Decoder)
%%%%%%%%%%%%%%%%%%%%%%%%%%%%%%%%%%%%%%%%%%%%%%%%%%%%%%%%%%%%%%%%%%%%%%%%%%%%%%%%%

% 输出嵌入层 (Output Embedding) - 右移
\pic[shift={(6,0)}] {RoundedRectangle={
    name=output_embed,
    caption=,
    fill=embedding,
    width=3,
    height=0.8,
    cornerradius=5pt,
    bordercolor=black!50,
    borderwidth=0.8pt
}};
\node[layer_text, font=\scriptsize\bfseries] at (output_embed-anchor) {输出嵌入（右移）};

% 位置编码 (Positional Encoding) - 解码器侧
\node[circle, draw=black!50, fill=gray!30, minimum size=0.8cm] (pos_enc_right) at (6, 1.5) {};
\node[layer_text] at (pos_enc_right) {PE};
\node[layer_text] at (5, 1.5) {位置编码};

% 加法符号
\node[circle, draw=black!50, fill=white, minimum size=0.6cm] (add2) at (7.5, 1.5) {};
\node[layer_text] at (add2) {+};

% 解码器第一层
% Masked Multi-Head Self-Attention
\pic[shift={(6,3)}] {RoundedRectangle={
    name=dec_mask_attn,
    caption=,
    fill=attention,
    width=3,
    height=0.8,
    cornerradius=5pt,
    bordercolor=black!50,
    borderwidth=0.8pt
}};
\node[layer_text, font=\scriptsize\bfseries] at (dec_mask_attn-anchor) {掩码多头自注意力};

% Add & Norm 1 (decoder)
\pic[shift={(6,4.2)}] {RoundedRectangle={
    name=dec_norm1,
    caption=,
    fill=norm,
    width=3,
    height=0.6,
    cornerradius=5pt,
    bordercolor=black!50,
    borderwidth=0.8pt
}};
\node[layer_text, font=\scriptsize\bfseries] at (dec_norm1-anchor) {相加 \& 层归一化};

% Cross-Attention (Multi-Head Attention)
\pic[shift={(6,5.4)}] {RoundedRectangle={
    name=cross_attn,
    caption=,
    fill=attention!70,
    width=3,
    height=0.8,
    cornerradius=5pt,
    bordercolor=black!50,
    borderwidth=0.8pt
}};
\node[layer_text, font=\scriptsize\bfseries] at (cross_attn-anchor) {多头注意力};

% Add & Norm 2 (decoder)
\pic[shift={(6,6.6)}] {RoundedRectangle={
    name=dec_norm2,
    caption=,
    fill=norm,
    width=3,
    height=0.6,
    cornerradius=5pt,
    bordercolor=black!50,
    borderwidth=0.8pt
}};
\node[layer_text, font=\scriptsize\bfseries] at (dec_norm2-anchor) {相加 \& 层归一化};

% Feed Forward Network (decoder)
\pic[shift={(6,7.8)}] {RoundedRectangle={
    name=dec_ff,
    caption=,
    fill=feedforward,
    width=3,
    height=0.8,
    cornerradius=5pt,
    bordercolor=black!50,
    borderwidth=0.8pt
}};
\node[layer_text] at (dec_ff-anchor) {前馈神经网络};

% Add & Norm 3 (decoder)
\pic[shift={(6,9)}] {RoundedRectangle={
    name=dec_norm3,
    caption=,
    fill=norm,
    width=3,
    height=0.6,
    cornerradius=5pt,
    bordercolor=black!50,
    borderwidth=0.8pt
}};
\node[layer_text, font=\scriptsize\bfseries] at (dec_norm3-anchor) {相加 \& 层归一化};

% Linear Layer
\pic[shift={(6.25,10.5)}] {RoundedRectangle={
    name=linear,
    caption=,
    fill=linear,
    width=2.5,
    height=0.6,
    cornerradius=5pt,
    bordercolor=black!50,
    borderwidth=0.8pt
}};
\node[layer_text] at (linear-anchor) {线性层};

% Softmax
\pic[shift={(6.25,11.7)}] {RoundedRectangle={
    name=softmax,
    caption=,
    fill=softmax,
    width=2.5,
    height=0.6,
    cornerradius=5pt,
    bordercolor=black!50,
    borderwidth=0.8pt
}};
\node[layer_text] at (softmax-anchor) {Softmax};

% 解码器背景框 - 只框住6个主要组件，使用透明的解码器颜色填充
\draw[thick, fill=decoder, fill opacity=0.3, draw=decoder!80!black, rounded corners=10pt] (5.8, 2.5) rectangle (9.2, 9.8);
\node[title_text, color=decoder!80!black, align=center] at (4.5, 8.1) {解码器\\(Decoder)};

%%%%%%%%%%%%%%%%%%%%%%%%%%%%%%%%%%%%%%%%%%%%%%%%%%%%%%%%%%%%%%%%%%%%%%%%%%%%%%%%%
% 连接线部分 (Connections) - 展示数据流向和残差连接
%%%%%%%%%%%%%%%%%%%%%%%%%%%%%%%%%%%%%%%%%%%%%%%%%%%%%%%%%%%%%%%%%%%%%%%%%%%%%%%%%

% 编码器内部连接 - 主要的前向传播路径
\draw[connect] (input_embed-north) -- (add1);
\draw[connect] (pos_enc_left) -- (add1);
\draw[connect] (add1) -- (enc_attn1-south);
\draw[connect] (enc_attn1-north) -- (enc_norm1-south);
\draw[connect] (enc_norm1-north) -- (enc_ff1-south);
\draw[connect] (enc_ff1-north) -- (enc_norm2-south);

% 编码器的残差连接（跳跃连接） - 根据Transformer架构图绘制正确的残差结构
% 【重要说明】残差连接的核心思想：将输入直接"跳过"中间的计算层，直接加到输出上
% 这样做的好处：1) 缓解梯度消失问题 2) 使网络更容易训练 3) 保持信息流动

% 第一个残差块：多头注意力的输入直接跳跃连接到Add&Norm1
% 【残差公式】output = MultiHeadAttention(input) + input，然后进行LayerNorm
% 定义多头注意力输入点（来自位置编码后的加法输出）
\path (add1) -- (enc_attn1-south) coordinate[pos=0.1] (enc_skip_start1);
% 定义Add&Norm1的侧面连接点
\path (enc_norm1-east) coordinate (enc_skip_end1);
% 绘制第一个残差连接：从多头注意力输入绕过多头注意力模块直接连接到Add&Norm1
% 修正箭头方向：箭头应该指向Add&Norm层
\draw[skip] (enc_skip_start1) -- ++(1.8,0) |- ($(enc_skip_end1) + (0.3,0)$) -- (enc_skip_end1);

% 第二个残差块：前馈网络的输入直接跳跃连接到Add&Norm2  
% 【残差公式】output = FeedForward(input) + input，然后进行LayerNorm
% 定义前馈网络输入点（来自第一个Add&Norm的输出）
\path (enc_norm1-north) -- (enc_ff1-south) coordinate[pos=0.1] (enc_skip_start2);
% 定义Add&Norm2的侧面连接点
\path (enc_norm2-east) coordinate (enc_skip_end2);
% 绘制第二个残差连接：从前馈网络输入绕过前馈网络模块直接连接到Add&Norm2
% 修正箭头方向：箭头应该指向Add&Norm层
\draw[skip] (enc_skip_start2) -- ++(1.8,0) |- ($(enc_skip_end2) + (0.3,0)$) -- (enc_skip_end2);

% 解码器内部连接 - 主要的前向传播路径
\draw[connect] (output_embed-north) -- (add2);
\draw[connect] (pos_enc_right) -- (add2);
\draw[connect] (add2) -- (dec_mask_attn-south);
\draw[connect] (dec_mask_attn-north) -- (dec_norm1-south);
\draw[connect] (dec_norm1-north) -- (cross_attn-south);
\draw[connect] (cross_attn-north) -- (dec_norm2-south);
\draw[connect] (dec_norm2-north) -- (dec_ff-south);
\draw[connect] (dec_ff-north) -- (dec_norm3-south);
\draw[connect] (dec_norm3-north) -- (linear-south);
\draw[connect] (linear-north) -- (softmax-south);

% 解码器的残差连接（跳跃连接） - 根据Transformer架构图绘制正确的残差结构
% 【重要说明】解码器有3个残差连接，每个都遵循相同的模式：input + SubLayer(input)

% 第一个残差块：掩码多头注意力的输入直接跳跃连接到Add&Norm1
% 【残差公式】output = MaskedMultiHeadAttention(input) + input，然后进行LayerNorm
% 【掩码作用】确保当前位置只能看到之前的位置，防止"未来信息泄漏"
% 定义掩码多头注意力输入点（来自位置编码后的加法输出）
\path (add2) -- (dec_mask_attn-south) coordinate[pos=0.1] (dec_skip_start1);
% 定义Add&Norm1的侧面连接点
\path (dec_norm1-east) coordinate (dec_skip_end1);
% 绘制第一个残差连接：从掩码多头注意力输入绕过掩码多头注意力模块直接连接到Add&Norm1
% 修正箭头方向：箭头应该指向Add&Norm层
\draw[skip] (dec_skip_start1) -- ++(1.8,0) |- ($(dec_skip_end1) + (0.3,0)$) -- (dec_skip_end1);

% 第二个残差块：多头注意力（交叉注意力）的输入直接跳跃连接到Add&Norm2
% 【残差公式】output = MultiHeadAttention(Q=dec_input, K=enc_output, V=enc_output) + dec_input
% 【交叉注意力】Q来自解码器，K和V来自编码器，实现编码器-解码器信息交互
% 定义多头注意力输入点（来自第一个Add&Norm的输出）
\path (dec_norm1-north) -- (cross_attn-south) coordinate[pos=0.1] (dec_skip_start2);
% 定义Add&Norm2的侧面连接点
\path (dec_norm2-east) coordinate (dec_skip_end2);
% 绘制第二个残差连接：从多头注意力输入绕过多头注意力模块直接连接到Add&Norm2
% 修正箭头方向：箭头应该指向Add&Norm层
\draw[skip] (dec_skip_start2) -- ++(1.8,0) |- ($(dec_skip_end2) + (0.3,0)$) -- (dec_skip_end2);

% 第三个残差块：前馈网络的输入直接跳跃连接到Add&Norm3
% 【残差公式】output = FeedForward(input) + input，然后进行LayerNorm
% 【前馈网络】通常是两个线性变换加ReLU：FFN(x) = max(0, xW1 + b1)W2 + b2
% 定义前馈网络输入点（来自第二个Add&Norm的输出）
\path (dec_norm2-north) -- (dec_ff-south) coordinate[pos=0.1] (dec_skip_start3);
% 定义Add&Norm3的侧面连接点
\path (dec_norm3-east) coordinate (dec_skip_end3);
% 绘制第三个残差连接：从前馈网络输入绕过前馈网络模块直接连接到Add&Norm3
% 修正箭头方向：箭头应该指向Add&Norm层
\draw[skip] (dec_skip_start3) -- ++(1.8,0) |- ($(dec_skip_end3) + (0.3,0)$) -- (dec_skip_end3);

% 编码器到解码器的连接（交叉注意力的K,V输入）
% 【关键连接】编码器的输出作为解码器交叉注意力的Key和Value
% 这是Transformer实现序列到序列转换的核心机制
\draw[connect, blue!70, very thick] (enc_norm2-east) -- ($(enc_norm2-east) + (1.5,0)$) |- (cross_attn-west);
% 添加K,V标签 - 表示Key和Value来自编码器
\node[layer_text, color=blue!70, font=\tiny] at ($(enc_norm2-east) + (0.8,-0.3)$) {K,V};

% 输入输出标签 - 明确数据流的起点和终点
\node[layer_text, below=0.3cm of input_embed-south] {输入序列};
\node[layer_text, above=0.3cm of softmax-north] {输出概率};

% 添加"N×"标记表示多层堆叠 - 调整位置
% 【说明】实际应用中编码器和解码器都是多层堆叠的，通常N=6或更多
\node[layer_text, color=red!70] at (-0.5, 5) {N×};
\node[layer_text, color=red!70] at (5.5, 6) {N×};

% 添加图例说明 - 采用左右布局避免重叠
\node[title_text] at (4, 13) {Transformer 架构图};

% 左侧：基本图例说明
\node[layer_text, align=left] at (0.5, -3) {
    \textbf{图例说明：}\\
    \textcolor{red!70}{N×}: 表示该结构重复N层\\
    \textcolor{blue!70}{蓝色粗线}: 编码器向解码器传递K,V\\
    \textcolor{gray!60}{虚线}: 残差连接（跳跃连接）\\
    \textcolor{blue!70}{残差连接作用：}缓解梯度消失问题，便于训练深层网络
};

% 右侧：详细的残差连接说明
\node[layer_text, align=left, font=\tiny] at (7.5, -3.2) {
    \textbf{残差连接详解：}\\
    \textcolor{gray!60}{编码器残差连接：}\\
    • 第1个：多头注意力输入 → Add\&Norm1（跳过多头注意力计算）\\
    • 第2个：前馈网络输入 → Add\&Norm2（跳过前馈网络计算）\\
    \textcolor{gray!60}{解码器残差连接：}\\
    • 第1个：掩码注意力输入 → Add\&Norm1（跳过掩码注意力计算）\\
    • 第2个：交叉注意力输入 → Add\&Norm2（跳过交叉注意力计算）\\
    • 第3个：前馈网络输入 → Add\&Norm3（跳过前馈网络计算）
};

\end{tikzpicture}
\end{document}
